		在 \texttt{segtree} 的幺半群 $(S,op,e)$ 上再定义作用幺半群 $(F,composition,id)$，即可支持区间修改、区间查询和树上二分，单次操作均为 $O(\log n)$。

		\textbf{五件套约束}：\texttt{op} 必须满足结合律；\texttt{e} 是其单位元；\texttt{mapping(f,x)} 表示把作用 $f$ 施加到 $x$；\texttt{composition(f,g)} 表示“先 $g$ 后 $f$”；\texttt{id} 是恒等作用。作用还必须满足对 \texttt{op} 的分配律。

		\textbf{最大的坑}：\texttt{mapping} 拿不到区间边界。凡是更新公式需要长度，都要把 \texttt{len} 存进 $S$，让它随 \texttt{op} 一起合并。

		\textbf{接口}：除 \texttt{segtree} 的 \texttt{set/get/prod/all\_prod/max\_right/min\_left} 外，新增 \texttt{apply(p,f)} 与闭区间 \texttt{apply(l,r,f)}。因为查询前可能需要下推标记，\texttt{get/prod} 与两个二分接口不再是 \texttt{const}。

		\subsubsection{模板本体}

		\begin{minted}{cpp}
/* ACL atcoder/lazysegtree 自包含移植 + 闭区间
 * lazy_segtree(n) 合法位置 0..n; 区间 [l, r] 左闭右闭
 * 1-based 只碰 1..n; 0-based 只碰 0..n-1
 * 从 vector 建树时 v[i] 对应位置 i, 长度由 vector 决定
 * 空区间写 r == l - 1, prod 返回 e(), apply 直接返回
 */
template<class S, S (*op)(S, S), S (*e)(),
         class F, S (*mapping)(F, S),
         F (*composition)(F, F), F (*id)()>
struct lazy_segtree {
public:
    lazy_segtree() : lazy_segtree(0) {}

    explicit lazy_segtree(int n)
        : lazy_segtree(vector<S>(n + 1, e())) {}

    explicit lazy_segtree(const vector<S> &v)
        : _n(int(v.size())) {
        size = (int) bit_ceil((unsigned) _n);
        log = countr_zero((unsigned) size);
        d = vector<S>(2 * size, e());
        lz = vector<F>(size, id());
        for (int i = 0; i < _n; i++)
            d[size + i] = v[i];
        for (int i = size - 1; i >= 1; i--)
            update(i);
    }

    void set(int p, S x) {
        assert(0 <= p && p < _n);
        p += size;
        for (int i = log; i >= 1; i--) push(p >> i);
        d[p] = x;
        for (int i = 1; i <= log; i++) update(p >> i);
    }

    S get(int p) {
        assert(0 <= p && p < _n);
        p += size;
        for (int i = log; i >= 1; i--) push(p >> i);
        return d[p];
    }

    S prod(int l, int r) {
        r++;
        assert(0 <= l && l <= r && r <= _n);
        if (l == r) return e();
        l += size, r += size;
        for (int i = log; i >= 1; i--) {
            if (((l >> i) << i) != l) push(l >> i);
            if (((r >> i) << i) != r)
                push((r - 1) >> i);
        }
        S sml = e(), smr = e();
        while (l < r) {
            if (l & 1) sml = op(sml, d[l++]);
            if (r & 1) smr = op(d[--r], smr);
            l >>= 1, r >>= 1;
        }
        return op(sml, smr);
    }

    S all_prod() { return d[1]; }

    void apply(int p, F f) {
        assert(0 <= p && p < _n);
        p += size;
        for (int i = log; i >= 1; i--) push(p >> i);
        d[p] = mapping(f, d[p]);
        for (int i = 1; i <= log; i++) update(p >> i);
    }

    void apply(int l, int r, F f) {
        r++;
        assert(0 <= l && l <= r && r <= _n);
        if (l == r) return;
        l += size, r += size;
        for (int i = log; i >= 1; i--) {
            if (((l >> i) << i) != l) push(l >> i);
            if (((r >> i) << i) != r)
                push((r - 1) >> i);
        }
        {
            int l2 = l, r2 = r;
            while (l < r) {
                if (l & 1) all_apply(l++, f);
                if (r & 1) all_apply(--r, f);
                l >>= 1, r >>= 1;
            }
            l = l2, r = r2;
        }
        for (int i = 1; i <= log; i++) {
            if (((l >> i) << i) != l) update(l >> i);
            if (((r >> i) << i) != r)
                update((r - 1) >> i);
        }
    }

    template<class G>
    int max_right(int l, G g) {
        assert(0 <= l && l <= _n);
        assert(g(e()));
        if (l == _n) return _n - 1;
        l += size;
        for (int i = log; i >= 1; i--) push(l >> i);
        S sm = e();
        do {
            while (l % 2 == 0) l >>= 1;
            if (!g(op(sm, d[l]))) {
                while (l < size) {
                    push(l);
                    l = 2 * l;
                    if (g(op(sm, d[l]))) {
                        sm = op(sm, d[l]);
                        l++;
                    }
                }
                return (l - size) - 1;
            }
            sm = op(sm, d[l]);
            l++;
        } while ((l & -l) != l);
        return _n - 1;
    }

    template<class G>
    int min_left(int r, G g) {
        r++;
        assert(0 <= r && r <= _n);
        assert(g(e()));
        if (r == 0) return 0;
        r += size;
        for (int i = log; i >= 1; i--)
            push((r - 1) >> i);
        S sm = e();
        do {
            r--;
            while (r > 1 && (r % 2)) r >>= 1;
            if (!g(op(d[r], sm))) {
                while (r < size) {
                    push(r);
                    r = 2 * r + 1;
                    if (g(op(d[r], sm))) {
                        sm = op(d[r], sm);
                        r--;
                    }
                }
                return r + 1 - size;
            }
            sm = op(d[r], sm);
        } while ((r & -r) != r);
        return 0;
    }

private:
    int _n, size, log;
    vector<S> d;
    vector<F> lz;

    static unsigned bit_ceil(unsigned n) {
        unsigned x = 1;
        while (x < n) x *= 2;
        return x;
    }

    static int countr_zero(unsigned n) {
        return __builtin_ctz(n);
    }

    void update(int k) {
        d[k] = op(d[2 * k], d[2 * k + 1]);
    }

    void all_apply(int k, F f) {
        d[k] = mapping(f, d[k]);
        if (k < size)
            lz[k] = composition(f, lz[k]);
    }

    void push(int k) {
        all_apply(2 * k, lz[k]);
        all_apply(2 * k + 1, lz[k]);
        lz[k] = id();
    }
};
		\end{minted}

		\begin{multicols*}{2}

			\subsubsection{示例：区间加 + 区间和}

			区间长度必须存进 $S$；空出来的 0 号位置长度为 $0$，因此不会影响 1-based 查询。

			\begin{minted}{cpp}
// ACL_EXAMPLE: lazy_range_add_sum
struct S {
    ll sum;
    int len;
};
using F = ll;

S op(S a, S b) {
    return {a.sum + b.sum, a.len + b.len};
}
S e() { return {0, 0}; }
S mapping(F f, S x) {
    return {x.sum + f * x.len, x.len};
}
F composition(F f, F g) { return f + g; }
F id() { return 0; }

using SegAddSum = lazy_segtree<S, op, e, F,
        mapping, composition, id>;

// vector<S> a(n + 1, e());
// FOR(i, 1, n) a[i] = {x[i], 1};
// SegAddSum seg(a);
// seg.apply(l, r, add);
// ll ans = seg.prod(l, r).sum;
			\end{minted}

			\subsubsection{示例：仿射修改 + 一至三次幂和}

			用 $x\leftarrow ax+b$ 统一表达乘法、加法与赋值：乘 $v$ 是 \texttt{\{v,0\}}，加 $v$ 是 \texttt{\{1,v\}}，赋值 $v$ 是 \texttt{\{0,v\}}。节点同时维护 $\sum x$、$\sum x^2$、$\sum x^3$。

			\begin{minted}{cpp}
// ACL_EXAMPLE: lazy_affine_moments
const ll MOD = 998244353;

ll mul_mod(ll a, ll b) {
    return (ll)((__int128)a * b % MOD);
}

struct S {
    ll s1, s2, s3;
    int len;
};
struct F {
    ll mul, add; // x <- mul * x + add
};

S op(S x, S y) {
    return {(x.s1 + y.s1) % MOD,
            (x.s2 + y.s2) % MOD,
            (x.s3 + y.s3) % MOD,
            x.len + y.len};
}
S e() { return {0, 0, 0, 0}; }

S mapping(F f, S x) {
    ll m = f.mul, a = f.add;
    ll m2 = mul_mod(m, m), m3 = mul_mod(m2, m);
    ll a2 = mul_mod(a, a), a3 = mul_mod(a2, a);
    ll old1 = x.s1, old2 = x.s2, old3 = x.s3;
    ll n = x.len % MOD;
    ll s1 = (mul_mod(m, old1) + mul_mod(a, n)) % MOD;
    ll s2 = (mul_mod(m2, old2)
           + mul_mod(2 * mul_mod(m, a) % MOD, old1)
           + mul_mod(a2, n)) % MOD;
    ll s3 = (mul_mod(m3, old3)
           + mul_mod(3 * mul_mod(m2, a) % MOD, old2)
           + mul_mod(3 * mul_mod(m, a2) % MOD, old1)
           + mul_mod(a3, n)) % MOD;
    return {s1, s2, s3, x.len};
}

F composition(F f, F g) {
    // 先 g 后 f
    return {mul_mod(f.mul, g.mul),
            (mul_mod(f.mul, g.add) + f.add) % MOD};
}
F id() { return {1, 0}; }

using SegMoments = lazy_segtree<S, op, e, F,
        mapping, composition, id>;

// add v:    seg.apply(l, r, {1, v});
// multiply: seg.apply(l, r, {v, 0});
// assign v: seg.apply(l, r, {0, v});
			\end{minted}

			\subsubsection{示例：赋值/加 + 和、最值与位置}

			这一份合并“扶苏的问题”和“最大值位置”板子。最大值并列时取最左位置；整段赋值后，最左最大值就是该段左端点。

			\begin{minted}{cpp}
// ACL_EXAMPLE: lazy_assign_add_stats
struct S {
    ll sum, mn, mx;
    int mx_pos, left, len;
};
struct F {
    bool has_assign;
    ll assign, add;
};

S op(S a, S b) {
    if (a.len == 0) return b;
    if (b.len == 0) return a;
    int pos = a.mx >= b.mx ? a.mx_pos : b.mx_pos;
    return {a.sum + b.sum, min(a.mn, b.mn),
            max(a.mx, b.mx), pos, a.left,
            a.len + b.len};
}
S e() { return {0, 0, 0, 0, 0, 0}; }

S mapping(F f, S x) {
    if (x.len == 0) return x;
    if (f.has_assign) {
        ll v = f.assign + f.add;
        return {v * x.len, v, v,
                x.left, x.left, x.len};
    }
    x.sum += f.add * x.len;
    x.mn += f.add;
    x.mx += f.add;
    return x;
}

F composition(F f, F g) {
    if (f.has_assign) return f;
    g.add += f.add;
    return g;
}
F id() { return {false, 0, 0}; }

using SegStats = lazy_segtree<S, op, e, F,
        mapping, composition, id>;

// leaf i: {x, x, x, i, i, 1}
// assign: seg.apply(l, r, {true, v, 0});
// add:    seg.apply(l, r, {false, 0, v});
			\end{minted}

			\subsubsection{示例：最小代价位置的权值和}

			区间加只平移最小值，不改变哪些位置达到最小值，因此对应权值和保持不变。常用于只统计“当前代价最小”转移点的 DP。

			\begin{minted}{cpp}
// ACL_EXAMPLE: lazy_min_weight
const ll MOD = 998244353;

struct S {
    ll mn, weight;
    bool empty;
};
using F = ll;

S op(S a, S b) {
    if (a.empty) return b;
    if (b.empty) return a;
    if (a.mn < b.mn) return a;
    if (b.mn < a.mn) return b;
    return {a.mn, (a.weight + b.weight) % MOD, false};
}
S e() { return {0, 0, true}; }
S mapping(F f, S x) {
    if (!x.empty) x.mn += f;
    return x;
}
F composition(F f, F g) { return f + g; }
F id() { return 0; }

using SegMinWeight = lazy_segtree<S, op, e, F,
        mapping, composition, id>;

// 真实叶子必须显式 set 成 {cost, weight, false}
// seg.apply(l, r, delta);
// S best = seg.prod(l, r);
			\end{minted}

			\subsubsection{示例：权值线段树求最小众数}

			离散化后，每个叶子代表一个权值。区间加维护一段权值的出现次数；合并时先取最高频次，频次相同取较小权值。

			\begin{minted}{cpp}
// ACL_EXAMPLE: lazy_value_mode
struct S {
    ll cnt, value;
    bool empty;
};
using F = ll;

S op(S a, S b) {
    if (a.empty) return b;
    if (b.empty) return a;
    if (a.cnt != b.cnt) return a.cnt > b.cnt ? a : b;
    return a.value < b.value ? a : b;
}
S e() { return {0, 0, true}; }
S mapping(F f, S x) {
    if (!x.empty) x.cnt += f;
    return x;
}
F composition(F f, F g) { return f + g; }
F id() { return 0; }

using SegMode = lazy_segtree<S, op, e, F,
        mapping, composition, id>;

// leaf i: {0, compressed_values[i], false}
// seg.apply(left_value_id, right_value_id, delta);
// auto [frequency, value, unused] = seg.all_prod();
			\end{minted}

		\end{multicols*}
